% Frame-title strip, evidence tags, text and table helpers, the fork-tree symbol.
% Section strip on every frame: the current section shows its number and name in teal, the others their number.
\newcommand{\currentact}{0}
\newcommand{\setact}[1]{\renewcommand{\currentact}{#1}}
\newcommand{\actitem}[2]{\ifnum\currentact=#1{\color{Teal}\bfseries #1\enspace #2}\else{\color{Muted!65}#1}\fi}
% Evidence tag: \ev{\tagM} inside a frame title prints the tag on the part-label line.
\newcommand{\currentev}{}
\newcommand{\currentsurf}{}
\newcommand{\surf}[1]{\gdef\currentsurf{#1}}
\newcommand{\ev}[1]{\gdef\currentev{#1}}
\newcommand{\evtag}[2]{\tikz[baseline=(x.base)]\node[draw=#1,text=#1,line width=.6pt,rounded corners=1.5pt,inner xsep=2.4pt,inner ysep=1.3pt,font=\fontsize{7.6}{8.4}\selectfont\bfseries](x){#2};}
\newcommand{\tagM}{\evtag{Teal}{MEASURED}}
\newcommand{\tagB}{\evtag{Ink}{BUILT}}
\newcommand{\tagP}{\evtag{Gold}{DESIGN}}
\newcommand{\tagO}{\evtag{Muted}{PUBLISHED}}
\newcommand{\tagA}{\evtag{Gold!55!Muted}{ILLUSTRATION}}
\newif\ifbackup
\def\endlabel{REFERENCES}\newcommand{\setendlabel}[1]{\gdef\endlabel{#1}}
\newcommand{\actstrip}{\ifbackup{\color{Teal}\bfseries\endlabel}\else{\color{Teal}\bfseries\currentact\enspace\ifcase\currentact\or MOTIVATION\or EXAMPLES\or RUNTIME\or EVIDENCE\or ACCEPTANCE\or CONCLUSION\fi}\fi}
\setbeamertemplate{frametitle}{\gdef\currentev{}\gdef\currentsurf{}\sbox0{\insertframetitle}\vspace{.1cm}{\fontsize{8}{9.5}\selectfont\actstrip\ifx\currentsurf\empty\else\quad{\color{Gold}\bfseries\MakeUppercase{\currentsurf}}\fi\hfill\currentev\par}\vspace{.08cm}\insertframetitle\par\vspace{.15cm}}
\setbeamertemplate{footline}{\hspace{.8cm}{\fontsize{7}{8}\selectfont\color{Muted}Yossi Eliaz\quad Cambridge SRG\quad 15 October 2026}\hfill{\fontsize{8}{9}\selectfont\color{Muted}\ifbackup\else\insertframenumber\,/\,\inserttotalframenumber\fi}\hspace{.8cm}\vspace{.17cm}}
\setlength{\tabcolsep}{5pt}
\renewcommand{\arraystretch}{1.15}
\newcolumntype{Y}{>{\raggedright\arraybackslash}X}
\newcolumntype{L}[1]{>{\raggedright\arraybackslash}p{#1}}
\newcommand{\lead}[1]{{\fontsize{12.5}{15}\selectfont #1\par}\vspace{.2cm}}
\newcommand{\claim}[1]{\par\vspace{.2cm}{\fontsize{11.5}{14}\selectfont\color{Teal}\bfseries #1\par}}
\newcommand{\sources}[1]{\par\vfill{\fontsize{7.8}{9.2}\selectfont\color{Muted}#1\par}\vspace{.05cm}}
\newcommand{\tagline}[1]{\par{\fontsize{9}{11}\selectfont\color{Muted}#1\par}}
\newcommand{\term}[1]{\textbf{#1}}
\newcommand{\no}{{\color{Rust}\bfseries$\times$}}
\newenvironment{ticks}{\setlength{\leftmargini}{1.3em}\begin{itemize}\setlength{\itemsep}{1.5pt}}{\end{itemize}}
\newcommand{\yes}{{\color{Teal}\bfseries$\checkmark$}}
% Recurring symbol: the title slide's fork tree, inline beside the refrain (slides 10, 36, 50).
\newcommand{\forktree}[1][Teal]{\tikz[baseline={(0,-.55ex)},x=1em,y=1em]{\draw[#1,line width=.75pt] (0,0)--(.85,0)--(1.55,.42) (.85,0)--(1.55,0) (.85,0)--(1.55,-.42);\foreach \fx/\fy in {0/0,.85/0,1.55/.42,1.55/0,1.55/-.42}{\fill[#1] (\fx,\fy) circle(.12);}}}

% Secondary details are retained in source and extracted into the presenter guide.
\newcommand{\qadetail}[1]{}
% Recurring proposed continuation lifecycle; focus: 1 state, 2 evidence, 3 authority, 0 all.
\newcommand{\lifecycle}[1]{%
\def\ls{Ink!35}\def\le{Ink!35}\def\la{Ink!35}%
\ifnum#1=0\relax\def\ls{Teal}\def\le{Teal}\def\la{Teal}\fi%
\ifnum#1=1\relax\def\ls{Teal}\fi%
\ifnum#1=2\relax\def\le{Teal}\fi%
\ifnum#1=3\relax\def\la{Teal}\fi%
\begin{center}
\begin{tikzpicture}[x=1cm,y=1cm, every node/.style={align=center,font=\fontsize{9.5}{11}\selectfont}, lc/.style={draw=Ink!35,fill=Pale,text width=3.15cm,minimum height=.72cm,inner sep=3pt}, active/.style={draw=Teal,line width=1pt}]
\node[lc,draw=\ls] (a) at (0,0) {Identify useful state};
\node[lc,draw=\ls] (b) at (4,0) {Create continuations};
\node[lc,draw=\le] (c) at (8,0) {Return artifact + observations};
\node[lc,draw=\le] (d) at (8,-1.15) {Select or compose};
\node[lc,draw=\le] (e) at (4,-1.15) {Verify exact\\artifact};
\node[lc,draw=\la] (f) at (0,-1.15) {Controller\\authorizes\\publication};
\draw[flow] (a)--(b);\draw[flow] (b)--(c);\draw[flow] (c)--(d);\draw[flow] (d)--(e);\draw[flow] (e)--(f);
\end{tikzpicture}
\end{center}}
